% Options for packages loaded elsewhere
\PassOptionsToPackage{unicode}{hyperref}
\PassOptionsToPackage{hyphens}{url}
\PassOptionsToPackage{dvipsnames,svgnames*,x11names*}{xcolor}
%
\documentclass[
  11pt,
]{article}
\usepackage{lmodern}
\usepackage{amssymb,amsmath}
\usepackage{ifxetex,ifluatex}
\ifnum 0\ifxetex 1\fi\ifluatex 1\fi=0 % if pdftex
  \usepackage[T1]{fontenc}
  \usepackage[utf8]{inputenc}
  \usepackage{textcomp} % provide euro and other symbols
\else % if luatex or xetex
  \usepackage{unicode-math}
  \defaultfontfeatures{Scale=MatchLowercase}
  \defaultfontfeatures[\rmfamily]{Ligatures=TeX,Scale=1}
\fi
% Use upquote if available, for straight quotes in verbatim environments
\IfFileExists{upquote.sty}{\usepackage{upquote}}{}
\IfFileExists{microtype.sty}{% use microtype if available
  \usepackage[]{microtype}
  \UseMicrotypeSet[protrusion]{basicmath} % disable protrusion for tt fonts
}{}
\makeatletter
\@ifundefined{KOMAClassName}{% if non-KOMA class
  \IfFileExists{parskip.sty}{%
    \usepackage{parskip}
  }{% else
    \setlength{\parindent}{0pt}
    \setlength{\parskip}{6pt plus 2pt minus 1pt}}
}{% if KOMA class
  \KOMAoptions{parskip=half}}
\makeatother
\usepackage{xcolor}
\IfFileExists{xurl.sty}{\usepackage{xurl}}{} % add URL line breaks if available
\IfFileExists{bookmark.sty}{\usepackage{bookmark}}{\usepackage{hyperref}}
\hypersetup{
  colorlinks=true,
  linkcolor=blue,
  filecolor=Maroon,
  citecolor=Blue,
  urlcolor=blue,
  pdfcreator={LaTeX via pandoc}}
\urlstyle{same} % disable monospaced font for URLs
\usepackage[margin=1in]{geometry}
\usepackage{longtable,booktabs}
% Correct order of tables after \paragraph or \subparagraph
\usepackage{etoolbox}
\makeatletter
\patchcmd\longtable{\par}{\if@noskipsec\mbox{}\fi\par}{}{}
\makeatother
% Allow footnotes in longtable head/foot
\IfFileExists{footnotehyper.sty}{\usepackage{footnotehyper}}{\usepackage{footnote}}
\makesavenoteenv{longtable}
\setlength{\emergencystretch}{3em} % prevent overfull lines
\providecommand{\tightlist}{%
  \setlength{\itemsep}{0pt}\setlength{\parskip}{0pt}}
\setcounter{secnumdepth}{-\maxdimen} % remove section numbering

\author{}
\date{}

\begin{document}

\hypertarget{calibration-free-fp8-and-satotate-fp10-quantization-for-transformers-a-shannon-prime-experiment-design}{%
\section{Calibration-Free fp8 and Sato--Tate fp10 Quantization for
Transformers: A Shannon-Prime Experiment
Design}\label{calibration-free-fp8-and-satotate-fp10-quantization-for-transformers-a-shannon-prime-experiment-design}}

\textbf{A. Knack} (Shannon-Prime Project) \textbf{Companion systems
paper. Draft v0.3 --- 2026-05-16}

\textbf{v0.3 changelog.} v0.2 erroneously claimed \(p = 11\) was a split
prime in \(K = \mathbb{Q}(\sqrt{-163})\) with \(a_{11} = -6\) ordinary
reduction. The Shannon-Prime Test Suite (suite version 0.1, test
\texttt{PAPER-D-FIX}) caught this: \((-163/11) = -1\), so 11 is
\emph{inert}, \(E_{11}\) is supersingular, and \(a_{11} = 0\) exactly.
The smallest split prime in \(K\) is \(p = 41\) --- exactly the first
value of Euler's polynomial \(n^2 + n + 41\), which is no accident,
since \(h(-163) = 1\) is what makes that polynomial prime-rich.
Throughout v0.3 we replace \(p_2 = 11\) with \(p_2 = 41\),
\(a_{p_2} = 1\), and the Frobenius polynomial becomes
\(\varphi_{41}^2 - \varphi_{41} + 41 = 0\).

\begin{center}\rule{0.5\linewidth}{0.5pt}\end{center}

\hypertarget{abstract}{%
\subsection{Abstract}\label{abstract}}

This companion paper specifies the experiment that probes the two
highest-leverage predictions of the Shannon-Prime framework: that (i)
fp8 quantization on a CM-encoded transformer state requires no
calibration, and (ii) the Sato--Tate inert-prime + split-prime
asymmetric construction yields a calibration-free fp10 mixed-precision
format with zero variance on the inert channel. The predictions are
direct consequences of Theorems 2 and 3 of the theory companion
{[}SP-Theory-Companion 2026{]}. We present the experimental design on a
Phi-3 reference model, with the implementation hooks already present in
the Shannon-Prime engine. We describe five configurations (baseline
fp16, SP-fp8 calibration-free, GPTQ-fp8 calibrated, AWQ-fp8 calibrated,
SP-fp10 Sato--Tate asymmetric).

\begin{center}\rule{0.5\linewidth}{0.5pt}\end{center}

\hypertarget{the-question}{%
\subsection{1. The Question}\label{the-question}}

Standard transformer quantization treats the model's weight and
activation tensors as floating-point arrays and rounds them to a
lower-precision representation. The rounding introduces errors in the
multiplicative relations the model has learned during training, and
post-training calibration is required to fit a per-tensor scale factor
that approximately restores those relations. GPTQ, AWQ, SmoothQuant, and
QServe all follow this pattern.

Shannon-Prime predicts something different. Its theoretical content
(Theorems 2 and 3 of {[}SP-Theory-Companion 2026{]}) is that on a CM
elliptic curve over \(\mathbb{Q}(\sqrt{-163})\), the Frobenius
endomorphism \(\varphi_p\) commutes with the layer endomorphisms. If the
model's state is encoded on this curve, then reduction to lower
precision is Frobenius application rather than rounding, and the
multiplicative relations are preserved exactly. Furthermore, the
inert/split prime partition under CM Sato--Tate gives an asymmetric
mixed-precision format with provably zero drift on the inert channel.

\textbf{Question.} \emph{Does SP-fp8 on a Phi-3 model recover
vanilla-fp16 perplexity to within noise, with no calibration data? Does
SP-fp10 Sato--Tate (2-bit inert + approximately 8-bit split at
\(p_2 = 41\)) match or beat SP-fp8 with no calibration?}

\begin{center}\rule{0.5\linewidth}{0.5pt}\end{center}

\hypertarget{setup}{%
\subsection{2. Setup}\label{setup}}

\hypertarget{reference-model}{%
\subsubsection{2.1 Reference Model}\label{reference-model}}

Phi-3 was chosen for three reasons: it is the SP model-pack's calibrated
reference (\(\Delta\mathrm{PPL} = +2.44\%\)); it benchmarks quickly
(3.8B parameters fit on a single A100); and it does not exhibit the
Qwen3 edge-fail artifact.

\hypertarget{bench-corpora}{%
\subsubsection{2.2 Bench Corpora}\label{bench-corpora}}

Two corpora at three context lengths each:

\begin{longtable}[]{@{}lll@{}}
\toprule
Corpus & Description & Tokens\tabularnewline
\midrule
\endhead
\textbf{WikiText-103} & Standard perplexity bench & 245M\tabularnewline
\textbf{The Stack v2 (filtered)} & Code completion bench &
100M\tabularnewline
\bottomrule
\end{longtable}

Context lengths: 512, 2048, 8192.

\hypertarget{configurations}{%
\subsubsection{2.3 Configurations}\label{configurations}}

Five configurations, all evaluated on the same corpus + context
combinations. Note the corrected prime \(p_2 = 41\) (was incorrectly
\(p = 11\) in v0.2).

\begin{longtable}[]{@{}llll@{}}
\toprule
\begin{minipage}[b]{0.22\columnwidth}\raggedright
Config\strut
\end{minipage} & \begin{minipage}[b]{0.22\columnwidth}\raggedright
Quant\strut
\end{minipage} & \begin{minipage}[b]{0.22\columnwidth}\raggedright
Calibration\strut
\end{minipage} & \begin{minipage}[b]{0.22\columnwidth}\raggedright
Notes\strut
\end{minipage}\tabularnewline
\midrule
\endhead
\begin{minipage}[t]{0.22\columnwidth}\raggedright
\textbf{A. Baseline}\strut
\end{minipage} & \begin{minipage}[t]{0.22\columnwidth}\raggedright
fp16\strut
\end{minipage} & \begin{minipage}[t]{0.22\columnwidth}\raggedright
---\strut
\end{minipage} & \begin{minipage}[t]{0.22\columnwidth}\raggedright
Ground truth\strut
\end{minipage}\tabularnewline
\begin{minipage}[t]{0.22\columnwidth}\raggedright
\textbf{B. SP-fp8 calibration-free}\strut
\end{minipage} & \begin{minipage}[t]{0.22\columnwidth}\raggedright
SP-fp8 (\(p = 41\) split, \(\varphi_p^8\))\strut
\end{minipage} & \begin{minipage}[t]{0.22\columnwidth}\raggedright
None\strut
\end{minipage} & \begin{minipage}[t]{0.22\columnwidth}\raggedright
Primary hypothesis (single-prime Frobenius, smallest split)\strut
\end{minipage}\tabularnewline
\begin{minipage}[t]{0.22\columnwidth}\raggedright
\textbf{C. GPTQ-fp8}\strut
\end{minipage} & \begin{minipage}[t]{0.22\columnwidth}\raggedright
fp8\strut
\end{minipage} & \begin{minipage}[t]{0.22\columnwidth}\raggedright
approximately 1000 samples\strut
\end{minipage} & \begin{minipage}[t]{0.22\columnwidth}\raggedright
Industry baseline\strut
\end{minipage}\tabularnewline
\begin{minipage}[t]{0.22\columnwidth}\raggedright
\textbf{D. AWQ-fp8}\strut
\end{minipage} & \begin{minipage}[t]{0.22\columnwidth}\raggedright
fp8\strut
\end{minipage} & \begin{minipage}[t]{0.22\columnwidth}\raggedright
approximately 1000 samples\strut
\end{minipage} & \begin{minipage}[t]{0.22\columnwidth}\raggedright
Industry baseline\strut
\end{minipage}\tabularnewline
\begin{minipage}[t]{0.22\columnwidth}\raggedright
\textbf{E. SP-fp10 Sato--Tate asymmetric}\strut
\end{minipage} & \begin{minipage}[t]{0.22\columnwidth}\raggedright
\(\varphi_{p_1}^{k_1} \circ \varphi_{p_2}^{k_2}\), \(p_1 = 2\) inert,
\(p_2 = 41\) split\strut
\end{minipage} & \begin{minipage}[t]{0.22\columnwidth}\raggedright
None\strut
\end{minipage} & \begin{minipage}[t]{0.22\columnwidth}\raggedright
Sharper hypothesis: zero-drift mixed precision\strut
\end{minipage}\tabularnewline
\bottomrule
\end{longtable}

The hypotheses are: - B matches A within noise:
\(|\Delta\mathrm{PPL}_{B,A}| < 0.5\%\). - B matches C/D within noise:
\(|\Delta\mathrm{PPL}_{B,C}|, |\Delta\mathrm{PPL}_{B,D}| < 0.3\%\). - E
at 10-bit total matches B at 8-bit total:
\(|\Delta\mathrm{PPL}_{E,B}| < 0.2\%\) despite lower precision, because
the inert channel contributes zero drift.

\begin{center}\rule{0.5\linewidth}{0.5pt}\end{center}

\hypertarget{why-fp8-why-p-41-why-varphi_418}{%
\subsection{\texorpdfstring{3. Why fp8, Why \(p = 41\), Why
\(\varphi_{41}^8\)}{3. Why fp8, Why p = 41, Why \textbackslash varphi\_\{41\}\^{}8}}\label{why-fp8-why-p-41-why-varphi_418}}

The Frobenius Quantization Theorem implements
\(\mathrm{fp}16 \to \mathrm{fp}q\) as \(\varphi_p^{16-q}\). For
\(q = 8\) at one byte of representation, \(p\) can be chosen up to
\(2^8 = 256\).

We choose \(p = 41\), the smallest split prime in \(K\). The Frobenius
element is \(\pi = \omega\) (or \(\bar{\omega}\); conjugate pair) with
\(N(\omega) = 41\). The Frobenius trace is
\(a_{41} = \mathrm{Tr}(\omega) = 1\), giving Frobenius characteristic
polynomial

\[\varphi_{41}^2 - \varphi_{41} + 41 = 0 \quad \text{in } \operatorname{End}(E_{41}).\]

By the Hasse--Weil decomposition,
\(\#E_{41}(\mathbb{F}_{41}) = 42 - 1 = 41\), well within the bound
\(41 + 1 + 2\sqrt{41} \approx 54.8\). The Sato--Tate angle is

\[\cos\theta_{41} = \frac{a_{41}}{2\sqrt{41}} = \frac{1}{2\sqrt{41}} \approx 0.0781,\]

which is \emph{unusually small} --- the corresponding drift per
Frobenius application is bounded by \(|a_{41}| = 1\), the smallest
possible non-zero magnitude for a split prime. Picking the smallest
split prime is therefore picking the smallest-drift split prime as well,
a happy alignment.

\textbf{Per-coordinate information capacity at \(q = 8\), \(p = 41\):}
\(\log_2(41 + 1 + 2\sqrt{41}) \approx 5.78\) bits, matching empirical
fp8 LLM weight bit-widths {[}Liu et al.~2024 KIVI{]}.

The choice \(p = 41\) also makes the Sato--Tate inert/split partition
concrete: \(\omega\) itself is the Frobenius element, and its conjugate
\(\bar{\omega} = 1 - \omega\) generates the conjugate ideal. The test
suite asserts the \emph{norm invariant}
\(N(\varphi_p^k(\text{state})) = N(\text{state}) \cdot p^k\) rather than
fixing a representative (since pi vs pi-bar is a choice).

\begin{center}\rule{0.5\linewidth}{0.5pt}\end{center}

\hypertarget{the-satotate-mixed-precision-choice-config-e}{%
\subsection{3.5 The Sato--Tate Mixed-Precision Choice (Config
E)}\label{the-satotate-mixed-precision-choice-config-e}}

Theorem 3 of the theory companion {[}SP-Theory-Companion 2026, §3A{]}
sharpens Theorem 2 by partitioning primes into \emph{inert}
(deterministic \(a_p = 0\), zero drift) and \emph{split} (bounded
analytic drift).

\textbf{Prime selection.} For \(K = \mathbb{Q}(\sqrt{-163})\): smallest
inert primes (Legendre \(-1\)) are
\(p_1 \in \{2, 3, 5, 7, 11, 13, 17, 19, 23, 29, 31, 37, \dots\}\);
smallest split primes (Legendre \(+1\)) are exactly the values
\(n^2 + n + 41\) that are prime:
\(\{41, 43, 47, 53, 61, 71, 83, 97, 113, 131, 151, 173, 197, \dots\}\).

Choose \(p_1 = 2\) (inert, \(\varphi_2^2 = -2\), scalar action with zero
\(\omega\)-drift) and \(p_2 = 41\) (split, smallest, with the smallest
split-prime drift \(|a_{41}| = 1\)).

\textbf{Resulting format.} Total bit-width approximately
\(2 + 8 \approx 10\) bits per encoded coordinate.

\begin{longtable}[]{@{}llll@{}}
\toprule
Format & Total bits & Inert (zero-drift) & Split
(bounded-drift)\tabularnewline
\midrule
\endhead
fp8 (Config B, \(p = 41\)) & 8 & 0 & 8\tabularnewline
fp10 Sato--Tate (Config E, \(p_1=2\), \(p_2=41\)) & 10 & 2 &
8\tabularnewline
fp16 (Config A) & 16 & 0 & 16\tabularnewline
\bottomrule
\end{longtable}

\textbf{Why Config E is informative.} Config B tests single-prime
Frobenius. Config E tests \emph{that} plus the inert/split partition. If
B passes and E fails, the inert-channel zero-drift prediction is wrong
despite Frobenius being correct --- a sharper signal than B alone.

\textbf{Suite-verified algebraic primitives.} The Shannon-Prime Test
Suite (v0.1, in
\texttt{D:\textbackslash{}F\textbackslash{}shannon-prime-repos\textbackslash{}test-suite\textbackslash{}})
verifies the following bit-exact properties used by Config E:

\begin{itemize}
\tightlist
\item
  \texttt{HOOK-B} (\texttt{-\/-frobenius-quant}):
  \(N(\varphi_{41}^k(\text{state})) = N(\text{state}) \cdot 41^k\) for
  arbitrary \(k\).
\item
  \texttt{HOOK-E} (\texttt{-\/-sato-tate-mix}):
  \(N(Q_{2,2,41,k_2}(\text{state})) = N(\text{state}) \cdot 4 \cdot 41^{k_2}\),
  and the inert channel applied alone scales state by \(-2\) exactly
  (zero \(\omega\)-drift).
\item
  \texttt{HOOK-E-COMM}: composition order is irrelevant --- applying
  \(\varphi_2^2\) then \(\varphi_{41}^k\) gives bit-identical output to
  applying \(\varphi_{41}^k\) then \(\varphi_2^2\) (50 random states
  tested).
\end{itemize}

These are pre-conditions for the A100 perplexity experiment; the
production engine kernels for \texttt{-\/-frobenius-quant} and
\texttt{-\/-sato-tate-mix} must reproduce the same bit-exact outputs.

\begin{center}\rule{0.5\linewidth}{0.5pt}\end{center}

\hypertarget{implementation}{%
\subsection{4. Implementation}\label{implementation}}

\hypertarget{existing-hooks}{%
\subsubsection{4.1 Existing Hooks}\label{existing-hooks}}

The Shannon-Prime engine already contains:

\begin{itemize}
\tightlist
\item
  \texttt{-\/-hier-ternary-mask} / \texttt{-\/-hier-res-bits-v} flags
  for ternary skeleton compression.
\item
  VHT2 + Möbius + spinor + square-free chain across all four backends.
\item
  Calibrated fp8 quantization in the engine path.
\item
  Per-architecture compression-default registry (model-pack scaffold);
  phi3 row exists and is calibrated.
\end{itemize}

What is needed:

\begin{enumerate}
\def\labelenumi{\arabic{enumi}.}
\tightlist
\item
  \textbf{A \texttt{-\/-frobenius-quant} flag} that toggles Frobenius
  reduction \(\varphi_{41}^k\) in place of calibrated fp8 (Config B).
\item
  \textbf{A \texttt{-\/-sato-tate-mix\ p1,k1,p2,k2} flag} that activates
  the Sato--Tate asymmetric mixed-precision encoding (Config E). For the
  production reference: \texttt{-\/-sato-tate-mix\ 2,2,41,8}.
\item
  \textbf{Unit tests} verifying:

  \begin{itemize}
  \tightlist
  \item
    \texttt{find\_element\_of\_norm(41)} returns an element of norm 41
    (i.e., \(\omega\) or \(\bar{\omega}\) --- either is fine, both give
    the same Frobenius action up to the well-known
    \(\pi \leftrightarrow \bar{\pi}\) ambiguity).
  \item
    \(N(\varphi_{41}^8(\text{state})) = N(\text{state}) \cdot 41^8\)
    exactly.
  \item
    \(\varphi_2^2\) acts as scalar \(-2\) (zero \(\omega\)-drift).
  \item
    \(\varphi_2^2 \circ \varphi_{41}^{k_2}\) matches
    \(\varphi_{41}^{k_2} \circ \varphi_2^2\) to bit-exactness, verifying
    commutativity (for E).
  \end{itemize}
\item
  \textbf{A bench harness modification} to disable calibration for B/E
  and to run five-seed variance measurement on E.
\end{enumerate}

The Python reference implementations at
\texttt{test-suite/src/engine\_hooks2.py} are bit-exact oracles; the
C/CUDA versions must agree.

Estimated implementation: 2--3 days at the engine level.

\hypertarget{outline-of-the-run-script}{%
\subsubsection{4.2 Outline of the Run
Script}\label{outline-of-the-run-script}}

\begin{verbatim}
for cfg in [A_baseline_fp16, B_sp_fp8_calfree_p41, C_gptq_fp8, D_awq_fp8, E_sp_fp10_satotate_2_41]:
    model = load_phi3(cfg.quantization)
    for corpus in [wikitext103, stackv2]:
        for ctx in [512, 2048, 8192]:
            ppl = eval_perplexity(model, corpus, ctx)
            tokens_per_sec = bench_throughput(model, corpus, ctx)
            log(cfg, corpus, ctx, ppl, tokens_per_sec)
\end{verbatim}

\hypertarget{passfail-criteria}{%
\subsubsection{4.3 Pass/Fail Criteria}\label{passfail-criteria}}

\begin{longtable}[]{@{}ll@{}}
\toprule
\begin{minipage}[b]{0.47\columnwidth}\raggedright
Comparison\strut
\end{minipage} & \begin{minipage}[b]{0.47\columnwidth}\raggedright
Criterion\strut
\end{minipage}\tabularnewline
\midrule
\endhead
\begin{minipage}[t]{0.47\columnwidth}\raggedright
\textbf{B vs A}\strut
\end{minipage} & \begin{minipage}[t]{0.47\columnwidth}\raggedright
\(|\Delta\mathrm{PPL}| < 0.5\%\) on WikiText-103 at all contexts\strut
\end{minipage}\tabularnewline
\begin{minipage}[t]{0.47\columnwidth}\raggedright
\textbf{B vs C}\strut
\end{minipage} & \begin{minipage}[t]{0.47\columnwidth}\raggedright
\(|\Delta\mathrm{PPL}_{B,C}| < 0.3\%\) at all contexts\strut
\end{minipage}\tabularnewline
\begin{minipage}[t]{0.47\columnwidth}\raggedright
\textbf{B vs D}\strut
\end{minipage} & \begin{minipage}[t]{0.47\columnwidth}\raggedright
\(|\Delta\mathrm{PPL}_{B,D}| < 0.3\%\) at all contexts\strut
\end{minipage}\tabularnewline
\begin{minipage}[t]{0.47\columnwidth}\raggedright
\textbf{B throughput}\strut
\end{minipage} & \begin{minipage}[t]{0.47\columnwidth}\raggedright
within 5\% of A throughput\strut
\end{minipage}\tabularnewline
\begin{minipage}[t]{0.47\columnwidth}\raggedright
\textbf{Robustness}\strut
\end{minipage} & \begin{minipage}[t]{0.47\columnwidth}\raggedright
B's relative PPL gap does not increase with context length\strut
\end{minipage}\tabularnewline
\begin{minipage}[t]{0.47\columnwidth}\raggedright
\textbf{E vs A}\strut
\end{minipage} & \begin{minipage}[t]{0.47\columnwidth}\raggedright
\(|\Delta\mathrm{PPL}_{E,A}| < 0.5\%\) at all contexts\strut
\end{minipage}\tabularnewline
\begin{minipage}[t]{0.47\columnwidth}\raggedright
\textbf{E vs B}\strut
\end{minipage} & \begin{minipage}[t]{0.47\columnwidth}\raggedright
\(|\Delta\mathrm{PPL}_{E,B}| < 0.2\%\) --- fp10 mixed precision matches
or beats fp8 single prime\strut
\end{minipage}\tabularnewline
\begin{minipage}[t]{0.47\columnwidth}\raggedright
\textbf{E zero-drift signature}\strut
\end{minipage} & \begin{minipage}[t]{0.47\columnwidth}\raggedright
Variance of \(\Delta\mathrm{PPL}\) across 5 random seeds is less than
\(0.5\times\) variance of B; verifies the inert-channel zero-drift
property\strut
\end{minipage}\tabularnewline
\bottomrule
\end{longtable}

If all seven criteria pass, both Theorem 2 and Theorem 3 are validated.
SP-fp8 ships as a calibration-free quantization; SP-fp10 ships as a
calibration-free \emph{variance-free} mixed precision.

\begin{center}\rule{0.5\linewidth}{0.5pt}\end{center}

\hypertarget{expected-outcome-and-risk-analysis}{%
\subsection{5. Expected Outcome and Risk
Analysis}\label{expected-outcome-and-risk-analysis}}

\hypertarget{expected-outcome}{%
\subsubsection{5.1 Expected Outcome}\label{expected-outcome}}

Theorem 2 predicts B matches A within noise. Expected magnitude of
\(|\Delta\mathrm{PPL}_{B,A}|\): 0.1\% to 0.5\%. With \(p = 41\) rather
than the erroneous \(p = 11\), the per-coordinate information capacity
is 5.78 bits rather than 4.17 bits --- Config B should be \emph{easier}
than the v0.2 paper predicted because there's more headroom.

Theorem 3 predicts E matches B within 0.2\%. The drift bound
\(|a_{41}| = 1\) at \(p_2 = 41\) is smaller than \(|a_{11}|\) would have
been \emph{if} 11 were split, so the corrected experiment has a sharper
prediction.

\hypertarget{what-a-negative-result-would-indicate}{%
\subsubsection{5.2 What a Negative Result Would
Indicate}\label{what-a-negative-result-would-indicate}}

A B-failure localizes to: (1) imperfect CM realization of Phi-3
endomorphisms (debuggable by mechanistic interpretability); (2)
Frobenius implementation bug (cross-check against the test-suite Python
oracle); (3) framework's curve choice doesn't match trained Phi-3.

An E-failure (with B passing) localizes specifically to the inert-prime
zero-drift prediction --- would falsify Lemma 3.1(a) for our model.

\hypertarget{what-a-positive-result-enables}{%
\subsubsection{5.3 What a Positive Result
Enables}\label{what-a-positive-result-enables}}

A positive result delivers: calibration-free fp8 deployment; a bridge to
fp4 (Theorem 2 at \(q = 4\)); calibration-free fp10 mixed precision
(Theorem 3); a target for ARM/Hexagon production.

\begin{center}\rule{0.5\linewidth}{0.5pt}\end{center}

\hypertarget{implementation-timeline}{%
\subsection{6. Implementation Timeline}\label{implementation-timeline}}

\begin{longtable}[]{@{}lll@{}}
\toprule
\begin{minipage}[b]{0.30\columnwidth}\raggedright
Phase\strut
\end{minipage} & \begin{minipage}[b]{0.30\columnwidth}\raggedright
Deliverable\strut
\end{minipage} & \begin{minipage}[b]{0.30\columnwidth}\raggedright
Estimated effort\strut
\end{minipage}\tabularnewline
\midrule
\endhead
\begin{minipage}[t]{0.30\columnwidth}\raggedright
0. \textbf{DONE}\strut
\end{minipage} & \begin{minipage}[t]{0.30\columnwidth}\raggedright
Shannon-Prime Test Suite v0.1 (algebraic oracle, 16 VERIFIED / 2 PENDING
/ 1 paper-flag)\strut
\end{minipage} & \begin{minipage}[t]{0.30\columnwidth}\raggedright
---\strut
\end{minipage}\tabularnewline
\begin{minipage}[t]{0.30\columnwidth}\raggedright
1. Engine hooks\strut
\end{minipage} & \begin{minipage}[t]{0.30\columnwidth}\raggedright
\texttt{-\/-frobenius-quant}, \texttt{-\/-sato-tate-mix} flags wired
against the Python oracle\strut
\end{minipage} & \begin{minipage}[t]{0.30\columnwidth}\raggedright
2--3 days\strut
\end{minipage}\tabularnewline
\begin{minipage}[t]{0.30\columnwidth}\raggedright
2. Bench harness\strut
\end{minipage} & \begin{minipage}[t]{0.30\columnwidth}\raggedright
Modified eval loop for the five configs, 5-seed variance measurement on
E\strut
\end{minipage} & \begin{minipage}[t]{0.30\columnwidth}\raggedright
1--2 days\strut
\end{minipage}\tabularnewline
\begin{minipage}[t]{0.30\columnwidth}\raggedright
3. Reference runs\strut
\end{minipage} & \begin{minipage}[t]{0.30\columnwidth}\raggedright
All five configs at all six (corpus × context) combinations\strut
\end{minipage} & \begin{minipage}[t]{0.30\columnwidth}\raggedright
1--2 days on A100\strut
\end{minipage}\tabularnewline
\begin{minipage}[t]{0.30\columnwidth}\raggedright
4. Analysis + writeup\strut
\end{minipage} & \begin{minipage}[t]{0.30\columnwidth}\raggedright
Result table, variance comparison for B and E\strut
\end{minipage} & \begin{minipage}[t]{0.30\columnwidth}\raggedright
2 days\strut
\end{minipage}\tabularnewline
\begin{minipage}[t]{0.30\columnwidth}\raggedright
5. Companion v0.4\strut
\end{minipage} & \begin{minipage}[t]{0.30\columnwidth}\raggedright
Incorporate results into this paper as §5 (Results)\strut
\end{minipage} & \begin{minipage}[t]{0.30\columnwidth}\raggedright
1 day\strut
\end{minipage}\tabularnewline
\bottomrule
\end{longtable}

Total: 8--10 days of engine + bench work after Phase 0.

\begin{center}\rule{0.5\linewidth}{0.5pt}\end{center}

\hypertarget{connection-to-the-broader-sp-framework}{%
\subsection{7. Connection to the Broader SP
Framework}\label{connection-to-the-broader-sp-framework}}

This experiment is one component of a larger program. The Shannon-Prime
framework predicts five additional implementations, each with its own
experimental signature: Stern--Brocot RoPE, Weil pairing attention,
Hecke-eigenform embeddings, L-function activation oracle, LLL KV write,
Iwasawa-tower depth stability, BSD analytic-rank training. Each is the
subject of a future SP companion paper.

The present experiment --- calibration-free fp8 and fp10 --- is first in
line because it is the smallest implementation step that probes the
strongest theoretical predictions, and now (post-v0.3) uses the
algebraically correct split prime.

\begin{center}\rule{0.5\linewidth}{0.5pt}\end{center}

\hypertarget{conclusion}{%
\subsection{8. Conclusion}\label{conclusion}}

Shannon-Prime's Frobenius Quantization Theorem and Sato--Tate
Mixed-Precision Theorem together predict that fp8 and fp10 quantization
on a CM-encoded transformer state require no calibration data, with the
fp10 format additionally exhibiting zero variance on its inert-prime
channel. The present companion paper (v0.3, corrected) has specified the
experimental design that tests both predictions on a Phi-3 reference
model. A positive result validates the framework's strongest practical
claims and opens a path to fp4, calibration-free mixed precision, and
edge-device deployment without calibration. The experiment is small
(approximately 1 week of engineering + benchmarking), the implementation
hooks are already in place across the SP engine and llama.cpp fork, the
algebraic correctness has been verified by the Shannon-Prime Test Suite,
and the framework's predicted outcomes are sharp and falsifiable.

\begin{center}\rule{0.5\linewidth}{0.5pt}\end{center}

\hypertarget{references}{%
\subsection{References}\label{references}}

\begin{itemize}
\tightlist
\item
  Shannon-Prime Test Suite v0.1,
  \texttt{D:\textbackslash{}F\textbackslash{}shannon-prime-repos\textbackslash{}test-suite\textbackslash{}}
  (this work).
\item
  Companion theory: ``Three Theorems for Shannon-Prime Compression:
  Hasse--Weil as the Shannon Limit, Frobenius as Quantization, and CM
  Sato--Tate Mixed Precision.'' 2026.
\item
  Full theory: ``Shannon-Prime: A CM-Elliptic-Curve Framework for
  Transformer Computation.'' 2026.
\item
  Full systems: ``Shannon-Prime: Provably Exact KV-Cache Compression and
  Per-Layer Acceleration on Standard Inference Stacks.'' 2026.
\item
  Liu, Z. \emph{et al.} ``KIVI: A Tuning-Free Asymmetric 2bit
  Quantization for KV Cache.'' 2024.
\item
  Frantar, E. \emph{et al.} ``GPTQ: Accurate Post-Training Quantization
  for Generative Pre-trained Transformers.'' 2023.
\item
  Lin, J. \emph{et al.} ``AWQ: Activation-aware Weight Quantization for
  LLM Compression and Acceleration.'' 2024.
\item
  Phi-3 Technical Report, Microsoft, 2024.
\end{itemize}

\end{document}
